一维粒子质量\(m\)，\(V(x)=-g[\delta(x-a)+\delta(x+a)]\)，\(g>0,a>0\)，\(\hbar\)保留。束缚态能量\(E=-\hbar^2\kappa^2/(2m)\)，\(\kappa>0\)。
(1) 由波函数连续和导数跳跃推导偶态、奇态的束缚能方程。
(2) 证明奇态出现需要满足最小势阱间距条件，偶态始终存在。
(3) 在\(a\kappa_0\gg1\)时求两态能量的指数劈裂，\(\kappa_0=mg/\hbar^2\)。
(4) 讨论\(a\to0\)时两delta融合的单势阱极限及态数变化。
